% A CHARACTER FROM AN UNDEFINED FAMILY IS COMPLAINED ABOUT ONCE.
%
% Reading a character whose family still holds \nullfont draws
%
%   ! \textfont 13 is undefined (character A).
%
% and then EMPTIES THE FIELD it was read from, so that a second reading
% of the same field finds nothing and says nothing. That matters for
% \mathaccent, which asks its nucleus twice -- once for the skew the
% accent slides by, once for the box the accent goes over. HSTeX read it
% twice and complained twice.
%
% The complaint names the size the reading wanted, so the same character
% draws \textfont, \scriptfont or \scriptscriptfont depending on where
% it stands.
\catcode`\{=1 \catcode`\}=2 \catcode`\$=3 \catcode`\^=7
\font\tenrm=cmr10 \font\teni=cmmi10 \font\tensy=cmsy10 \font\tenex=cmex10
\textfont0=\tenrm \scriptfont0=\tenrm \scriptscriptfont0=\tenrm
\textfont1=\teni \scriptfont1=\teni \scriptscriptfont1=\teni
\textfont2=\tensy \scriptfont2=\tensy \scriptscriptfont2=\tensy
\textfont3=\tenex \scriptfont3=\tenex \scriptscriptfont3=\tenex
\tenrm
\tracingonline=1 \hsize=200pt \parindent=0pt \hbadness=10000 \hfuzz=100pt
\message{[A a mathaccent over a group in an undefined family]}
\setbox0=\vbox{\noindent$B\mathaccent"161 {\fam13 A}C$\par}
\message{[B a bare character as the nucleus]}
\setbox0=\vbox{\noindent$B\fam13 \mathaccent"161 AC$\par}
\message{[C no accent, just the character]}
\setbox0=\vbox{\noindent$B{\fam13 A}C$\par}
\message{[done]}
\end
